\section{Correspondence Argument and Its Remaining Boundary}\label{app:proof}

This appendix separates three elementary results from the unfinished
process-to-rule simulation. The elementary results concern the generated
Tamarin systems themselves. They do not establish the complete
ProVerif correspondence statement of Section~\ref{sec:method:thm}.

\subsection{Observation and the Admitted Property Fragment}

Fix the protocol events mentioned in a property. Initialization,
\texttt{Misuse}, \texttt{Clone} and \texttt{Rollback} are administrative
actions and are hidden. A trace is a finite sequence of action sets;
empty sets are deleted. Repeated occurrences of the same reveal event
may be removed only when reveal events occur positively in correspondence
conclusions or negatively in reachability formulas. They must not occur
in equality premises or in positive reachability requirements. Otherwise,
deduplication can change a property's truth value.

For an equality property over event timestamps, the timestamp denotes a
rule instance, not an individual argument or a message transmission.
Replacing two sequential ProVerif events by one Tamarin action set can
therefore change timestamp equality. The strengthened proof fragment
requires at most one observed event in each protocol block
(\texttt{ObsOne}). It additionally requires reveal-use restrictions
(\texttt{LeakUse}) and that the identity guarding a ProVerif leak be
tuple-projectable from the initialization output (\texttt{LeakGuard}).
These conditions are necessary restrictions, not a certificate of the
remaining simulation. Properties with misuse exceptions are compared
directly on their own variants; hiding misuse actions does not preserve
a formula that explicitly names those actions.

\subsection{Elementary Tamarin Inclusions}

\noindent\textbf{Proposition A.1 (standard runs embed in misuse runs).}
After hiding administrative actions, every finite run of $T$ is a run
of $T_{\mathsf{cl}}$ and of $T_{\mathsf{rb}}$ with the same observations.

\emph{Proof.} Replay initialization, protocol, responder and leak rule
instances without firing a misuse rule. The Clone system has the same
premises and conclusions for those instances. In the Rollback system,
initialization and each update additionally create a persistent snapshot;
these facts enable rules but disable none. Since no instance emits
\texttt{Misuse}, the universal \texttt{OneMisuse} restriction holds
vacuously. Hiding administrative actions leaves the original sequence.
\hfill$\square$

\noindent\textbf{Proposition A.2 (persistent-state overapproximation).}
After the same projection, every finite run of $T_{\mathsf{cl}}$ or
$T_{\mathsf{rb}}$ can be reproduced by $T^{\omega}$.

\emph{Proof.} Induct over the given run. Maintain the invariant that
every available linear state value and every saved snapshot has a
corresponding persistent \texttt{St} fact in $T^{\omega}$; persistent
thread and verification facts agree, and attacker knowledge has the same
deductive closure. Initialization establishes the invariant. A protocol
step can use the corresponding persistent state and the same attacker
input; it emits the same protocol actions and outputs and retains the
successor persistently. Responders and leaks have the same non-state
premises and outputs. A Clone step copies an already represented state,
so it is matched by no step. A Rollback step restores a snapshot whose
value was produced at initialization or an earlier update, so it too is
matched by no step. Snapshot and misuse actions are hidden. These cases
preserve the invariant and the observation. The one-misuse restriction
only reduces the source runs and need not be imposed on the target.
This is an inclusion, not the reverse inclusion.\hfill$\square$

\noindent\textbf{Proposition A.3 (finite update bound).}
If the reachable state-update graph of the standard system is acyclic
and has longest path length $m$, every thread of $T$ performs at most
$m$ protocol updates. Thus $T_{\le n}=T$ for $n\ge m$.

\emph{Proof.} Initialization creates one linear state for a fresh
thread. Each protocol update consumes that state and creates exactly
one successor. Responders and leaks do not create state; the standard
system has no Clone or Rollback rule. Consequently the sequence of
states for a thread is a path from its initial vertex. Its length is
bounded by the longest path. The argument does not bound the number of
threads, responder steps or leak steps, and does not apply to the
misuse variants.\hfill$\square$

\subsection{Ordered and Unordered Correspondences}

\noindent\textbf{Proposition A.4 (prefix argument).}
Consider a prefix-closed execution set and a universal single-premise
correspondence whose conclusions are disjunctions of conjunctions of
positive events. Suppose that no rule instance contains both a premise
event and a conclusion event of that correspondence. Validity with all
conclusion events strictly earlier than the premise is equivalent to
validity without those ordering constraints.

\emph{Proof.} The ordered formula implies the unordered one. Conversely,
if an ordered correspondence fails at a premise occurrence, truncate the
run immediately after that occurrence. Prefix closure preserves
admissibility. No conclusion occurs simultaneously with the premise by
the hypothesis. In each disjunct at least one required event is absent
from the prefix; otherwise that disjunct would have held earlier. The
prefix therefore violates the unordered formula. This establishes the
contrapositive.\hfill$\square$

The generated \texttt{OneMisuse} restriction is prefix closed: deleting
a suffix cannot introduce two distinct misuse occurrences for one
identity. This proves the required closure for that restriction, not
for arbitrary restrictions supplied by a user. For example, a
restriction requiring an event eventually to occur cannot be admitted
by this argument.

\subsection{A Checker Counterexample and the Unfinished Simulation}

In a negative-control model we remove the public identity output and
all signing outputs from LM-OTS, and ask whether its reveal event is
reachable. The old report contains no FAIL. ProVerif proves the event
unreachable, whereas Tamarin finds it reachable: the former waits for
an attacker input equal to an unpublished fresh identity; the latter
can fire the leak from the persistent thread fact. This positive reveal
query is deliberately outside the strengthened property fragment. It
demonstrates that the old checker did not enforce either identity
availability or the admitted reveal-use condition. Version~0.1.1 rejects
it with \texttt{LeakGuard} and \texttt{LeakUse}; both existing signing
CIMs pass the added checks. No saved experimental model or verdict is
replaced by this software update.

The outstanding proof must construct a relation between concrete
ProVerif configurations and multiset-rewriting states, rather than
identifying a pending private-channel output with a state fact at all
times. A worker may consume that output before receiving an attacker
message; it then holds the state locally. If the input fails a guard,
the state is lost silently. A trace simulation may retain the
corresponding source state without scheduling it again, but this is a
weak, finite-trace argument, not a bisimulation of configurations or a
liveness theorem. A completed proof must cover these reserved and
abandoned states explicitly.

The remaining obligations include serializing unfinished blocks and
relating initialization table insertions to public output. They also
include fresh name allocation, multiple outputs and partial prefixes,
and one bounded Rollback request including its attacker knowledge. The operational
semantics and the Horn abstraction are separate levels: equality of
two observed clause sets does not prove equality of process trace sets
or justify a general impossibility claim for all verifier options.
Propositions A.1--A.4 and passing regression tests do not discharge these
remaining obligations. Transfer-dependent Achievement Levels remain
conditional pending a reviewed process-to-rule proof.
